% Fragment for the standalone six-point manuscript; no preamble.
% BEGIN UNIVERSAL FOUNDATION
\section{Universal signed matchings and the original presentation}
\label{sec:universal}

We distinguish a relation in an integral presentation from a relation that
holds only after a particular specialization. This distinction is essential:
dividing an integral relation by its content can remove precisely the torsion
that creates a cyclic Z-pair.

\begin{lemma}[Occurrence matching]\label{lem:occurrence-matching}
Let $(a_i)_{i=0}^5$ and $(b_i)_{i=0}^5$ be two lists in an abelian group,
including multiplicities. Their directed autocorrelations agree if and only
if there are a bijection $\sigma$ of their fifteen unordered edge
occurrences and signs $\epsilon_{ij}\in\{1,-1\}$ such that
\[
 a_j-a_i=\epsilon_{ij}(b_l-b_k),\qquad
 (k,l)=\sigma(i,j),\quad i<j, k<l.
\]
\end{lemma}
\begin{proof}
Each unordered edge with difference $g$ contributes the two directed
occurrences $g,-g$. Partition differences into orbits under $g\mapsto-g$.
For an orbit of size two, its autocorrelation coefficient at either element
is the number of unordered edges in that orbit. For a nonzero orbit of size
one it is twice that number. At zero it is six plus twice the number of
coincident off-diagonal unordered edges. Consequently equality of
autocorrelations gives equal numbers of edge occurrences in every orbit.
Match those occurrences bijectively within each orbit and choose a sign
that gives equality. If $g=-g$, both signs are available and neither may be
discarded. Conversely, a signed bijection equates the two directed
occurrences contributed by each edge, and the six diagonal zeros agree.
\end{proof}

Write $\mathcal E=\{(i,j):0\le i<j\le5\}$. In $\Z^{10}$ use the row vectors
$e_1,\ldots,e_5,f_1,\ldots,f_5$ as the standard basis, and set
$e_0=f_0=0$. Independently translate each input endpoint so that
$a_0=b_0=0$. A signed matching defines the $15\times10$ matrix with rows
\begin{equation}\label{eq:formal-matching-row}
 m_{ij}=e_j-e_i-\epsilon_{ij}(f_l-f_k),\qquad(k,l)=\sigma(i,j).
\end{equation}
The \emph{integral relation lattice}, the \emph{original presentation},
and its free rank are
\begin{equation}\label{eq:original-presentation}
 L_M=\sum_{(i,j)\in\mathcal E}\Z m_{ij},\qquad
 G_M=\Z^{10}/L_M,\qquad d=10-\rank_{\mathbb Q}M.
\end{equation}
We write $E_i,F_i$ for the classes of $e_i,f_i$ in $G_M$. The two labelled
lists $(E_i)$ and $(F_i)$ are called the universal configurations.

\begin{lemma}[Universal property and rank]\label{lem:universal-property}
The universal configurations are homometric as lists. Every realization
of the signed matching in an abelian group $K$ factors through a
homomorphism $G_M\to K$. Moreover $5\le\rank_{\mathbb Q}M\le10$, so
$0\le d\le5$. A realization by six-element sets rules out universal
within-endpoint collisions; a $\TI$-distinct realization rules out
universal $\TI$ equivalence.
\end{lemma}
\begin{proof}
Equation~\eqref{eq:formal-matching-row} gives a signed edge bijection in
$G_M$, so Lemma~\ref{lem:occurrence-matching} gives homometry. Sending
$e_i$ to $a_i$ and $f_i$ to $b_i$ defines a map from $\Z^{10}$ that kills
every row of $M$, hence factors through the quotient. The first five
columns of $M$ are the reduced incidence matrix of the complete graph
on six vertices; the rows for edges $(0,i)$ contain the five standard
basis vectors, giving rank five. There are ten columns. Finally a
universal collision or rigid equivalence remains true under every
homomorphism and therefore under the given realization.
\end{proof}

The rank in~\eqref{eq:original-presentation} belongs to the original
signed matching. A map from $G_M$ into a cylinder or a finite cyclic group
can impose additional relations and lower its rank; this never changes
which branch the original presentation occupies. In particular, we never
seek a nonzero map from the finite target $C_n$ to $\R$.

\subsection{The exact meaning of a lattice certificate}
For row matrices $M,N$ with ten columns, two integer matrices satisfying
\begin{equation}\label{eq:two-way-lattice-cert}
 AM=N,\qquad BN=M
\end{equation}
certify $L_M=L_N$. An equality of rational row spaces does not certify this
claim. The latter equality would identify only the free solutions and
could identify different torsion groups. We use primitive rational rows
only when enumerating height subspaces below.

To certify coordinates for $G_M$, supply integer matrices $U,V,D$ with
\begin{equation}\label{eq:smith-contract}
 UMV=D,\qquad \det U,\det V\in\{1,-1\},
\end{equation}
where $U$ is $15\times15$, $V$ is $10\times10$, and $D$ is rectangular
diagonal with nonzero entries $d_1,\ldots,d_r$ and all other entries zero.
The map sending a row vector $z$ to $zV$ induces an isomorphism
\[
 G_M\cong\bigoplus_{i=1}^{r}\Z/|d_i|\Z\ \oplus\ \Z^{10-r}.
\]
Factors of order one are omitted. Thus the coordinates of an old generator
are its row of $V$, reduced modulo $|d_i|$ in each torsion position. All
of~\eqref{eq:two-way-lattice-cert}--\eqref{eq:smith-contract} can be
checked by integer multiplication and determinants. Discovery by Hermite
or Smith algorithms is useful; their claimed normal forms are not accepted
without these identities. A collision certificate can instead be a row
$c$ with $cM=e_i-e_j$ (or $f_i-f_j$), $i\ne j$. A rigid-equivalence
certificate consists of the corresponding six integer row identities.

\subsection{The finite torsion bound and the finite branch}
\begin{lemma}\label{lem:bound}
The torsion subgroup $T_M$ of the original presentation $G_M$ has order
at most $135$.
\end{lemma}
\begin{proof}
First every square minor of a reduced oriented incidence matrix has
determinant $0,1$ or $-1$. In a chosen square submatrix each row has
at most two entries; if some row has at most one nonzero entry, expansion
reduces the assertion by induction. Otherwise every row has two opposite
entries, so the sum of its columns is zero and its determinant vanishes.
Permuting or changing signs of rows preserves this property. Both
five-column blocks of $M$ have it.

An $r\times r$ minor using $c$ columns of the first block has a Laplace
expansion with $\binom rc$ summands, each a product of two incidence
minors of absolute value at most one. For $r\le9$ this is at most
$\binom94=126$. If $r=10$, all ten columns occur. The chosen ten rows
correspond to ten distinct edges of the first $K_6$. A five-row
first-block determinant is nonzero precisely for a spanning tree:
a cycle gives dependence; five acyclic edges on six vertices form a
connected tree; and leaf elimination gives determinant $\pm1$ for a
tree. The complementary second-block determinant still has absolute
value at most one. Therefore
\[
 |\det M_R|\le\tau(K_6[R]),\qquad |R|=10,
\]
where $\tau$ counts spanning trees. No assertion that these summands
have a common sign is needed.

The bound on $\tau$ is an explicit small finite obligation with no
symmetry reduction. Enumerate every $R\subset\mathcal E$ of size ten
and every $J\subset\mathcal E$ of size five. Check connectedness and
acyclicity of the graph with edges $J$, and for every tree count its
containment $J\subset R$. Independently, form the $6\times6$ graph
Laplacian of $R$, remove one fixed row and column, and compute the exact
$5\times5$ determinant. The two answers agree because, for the reduced
incidence matrix $B_R$, Cauchy--Binet gives
\[
 \det(B_R^{\mathsf T}B_R)
   =\sum_{\substack{J\subset R\\|J|=5}}(\det B_J)^2
   =\tau(K_6[R]).
\]
This also proves the Laplacian method's counting interpretation, rather
than using an unproved counting shortcut.

The acceptance predicate requires every one of the $\binom{15}{10}=3003$
labelled $R$ and every one of the $\binom{15}{5}=3003$ possible $J$,
and equality of the two counts for each $R$. The complete histogram is:
\begin{center}\small
\begin{tabular}{r|rrrrrrr}
\toprule
$\tau$ & $0$ & $75$ & $96$ & $99$ & $100$ & $104$ & $111$\\
Number of $R$ & $6$ & $300$ & $90$ & $360$ & $90$ & $360$ & $360$\\
\midrule
$\tau$ & $114$ & $115$ & $120$ & $121$ & $128$ & $130$ & $135$\\
Number of $R$ & $360$ & $360$ & $180$ & $72$ & $45$ & $360$ & $60$\\
\bottomrule
\end{tabular}
\end{center}
The counts sum to $3003$, and the maximum is $135$. The project and
the independent review reproduce the full histogram by the two methods.
The incidence argument together with this exact finite obligation bounds
every nonzero maximal-rank minor by $135$.

Finally the ideal generated by the rank-$r$ minors is unchanged by
unimodular row and column operations, by Cauchy--Binet applied to those
operations and their integer inverses. In a diagonal presentation it
is generated by $\prod_{i=1}^r d_i$, whose absolute value is $|T_M|$.
Thus $|T_M|$ is the positive gcd of the nonzero maximal-rank minors.
It is no greater than any one of their absolute values and therefore
at most $135$.
\end{proof}

\begin{corollary}\label{cor:finite}
If $d=0$ and $G_M\to C_n$ realizes six-element endpoints, then, after
their independent anchors are removed, both endpoints lie in a cyclic
subgroup of order $q$ with $6\le q\le135$ and $q\mid n$. They are an
ordinary subgroup inflation of two homometric six-subsets of $C_q$.
\end{corollary}
\begin{proof}
Now $G_M=T_M$ is finite. Its image $H$ in $C_n$ is a cyclic subgroup,
so its order $q$ divides $n$ and $|T_M|$, and hence is at most $135$.
Six distinct endpoint coordinates force $q\ge6$. Identify $H$ with
$C_q$ by the injective map $t\mapsto(n/q)t$. The inverse on $H$
transports the signed matching and the six-element endpoints to $C_q$.
Restoring the separate translations of the original endpoints does not
change their vertices in the grammar graph.
\end{proof}

\subsection{The real solution space and a fixed congruence}
Let
\[
 V_M=\Hom(G_M,\R)=\{v\in\R^{10}:Mv=0\}.
\]
This real vector space has dimension $d$. Its points are the coordinates
of real homometric six-atom lists, even when distinct universal vertices
have identical real coordinates.

\begin{lemma}[Finite unions of subspaces]\label{lem:finite-subspaces}
Over an infinite field a finite-dimensional vector space cannot be covered
by finitely many proper linear subspaces.
\end{lemma}
\begin{proof}
For each proposed proper subspace choose a nonzero linear functional
vanishing on it. Their product is a nonzero polynomial in a basis of the
whole space. A nonzero polynomial cannot vanish at all points over an
infinite field: in one variable it has only finitely many roots; induction
on the number of variables applies this fact to its coefficient
polynomials. The product therefore cannot vanish throughout the space.
\end{proof}

For a permutation $\pi$ of six labels and $\eta\in\{1,-1\}$, define
\[
 T_{\pi,\eta}=\{v\in\R^{10}:
  b_i=\eta(a_{\pi(i)}-a_{\pi(0)})\text{ for every }i\}.
\]
These are linear subspaces. Since $a_0=b_0=0$, every real translation or
reflection equivalence of the two multisets has exactly this form for
some $\pi,\eta$. Labels with equal coordinates can be matched arbitrarily;
there are still only $2\cdot6!$ choices.

\begin{lemma}[Fixed free-quotient alignment]\label{lem:fixed-alignment}
If the two real multisets are congruent for every $v\in V_M$, then one
fixed $T_{\pi,\eta}$ contains $V_M$. After relabelling, reflecting and
reanchoring the second universal configuration, every real solution
satisfies $b_i=a_i$ for all $i$. These changes are invertible integral
changes of generators and preserve $G_M$ and $d$.
\end{lemma}
\begin{proof}
The hypothesis says that the finitely many subspaces
$V_M\cap T_{\pi,\eta}$ cover $V_M$. Lemma~\ref{lem:finite-subspaces}
implies that one is the entire space. Relabelling and reflection give
$b'_i=\eta(b_{\pi^{-1}(i)}-b_{\pi^{-1}(0)})=a_i$ on $V_M$.
A permutation followed by reanchoring transforms the five coordinates
of a labelled endpoint by an integer linear map. Undoing the permutation
and reanchoring is its integer inverse. Reflection multiplies that block
by $-1$. Hence the full ten-coordinate transformation is unimodular.
\end{proof}

The conclusion identifies the free projections, not the group elements
$E_i,F_i$. In a finitely generated abelian group an element is killed by
every real homomorphism precisely when it is torsion: write the group as
$\Z^d\oplus T$ and use its $d$ coordinate maps. Consequently the aligned
$F_i-E_i$ may be nonzero torsion. Imposing $F_i=E_i$ as \emph{integral}
relations would discard the cyclic phenomena under investigation.

Substitute $b_i=a_i$ only in the real equations. With six unanchored
heights, the resulting rows have the form
\begin{equation}\label{eq:height-rows}
 (u_j-u_i)-\epsilon(u_l-u_k),\qquad i<j, k<l,
\end{equation}
where $u_0,\ldots,u_5$ denote the six standard vectors. They have coordinate
sum zero. Let $S$ be their rational span. Removing the common translation
by setting $a_0=0$ gives a five-dimensional height space, and
\begin{equation}\label{eq:height-codimension}
 \dim_{\mathbb Q}S=5-d.
\end{equation}
Indeed every vector in $V_M$ lies on the aligned graph. Conversely any
anchored height vector satisfying~\eqref{eq:height-rows}, placed on that
graph, satisfies the original real matching equations and lies in $V_M$.
Thus the anchored kernel of $S$ is exactly $V_M$.
% END UNIVERSAL FOUNDATION

% BEGIN WEIGHTED LINE INPUT AND BF
\section{The weighted real-line input and the Bloom branch}
\label{sec:formal-bloom}
This section states explicitly the mathematical formula checked by the
line solvers. The result is the previously reviewed I/Iw component, tagged
\textbf{[PROVED], in-house, computer-assisted}. Its completeness remains
conditional on the saved exact UNSAT results and their audited encodings;
the elementary interface proof below does not remove that dependency.

\begin{lemma}[Weighted six-atom line theorem Iw]\label{lem:iw}
Two noncongruent homometric real multisets of total multiplicity six
have, after independent translations and reflections and possible
interchange, one of the following weakly ordered forms:
\begin{align}
 A_1&=(0,q-2p,q,2q-3p,3q-2p,3q-p),\notag\\
 B_1&=(0,q-p,q+p,2q+p,3q-2p,3q-p),
       &&p>0,\quad q\ge3p;\label{eq:weighted-form-one}\\
 A_2&=(0,q-p,q+p,3q-2p,3q-p,2q+p),\notag\\
 B_2&=(0,p,2q-p,q+2p,3q-p,2q+p),
       &&p>0,\quad 3p/2\le q<2p.\label{eq:weighted-form-two}
\end{align}
For six-element sets the two lower bounds are strict. For integer
coordinates the recovered $p,q$ are integers. Every displayed pair is
homometric and noncongruent.
\end{lemma}

\begin{proof}[Mathematical and computational dependency]
The normalization, exact occurrence-count encoding, coefficient-plane
exclusions and solver obligation are given in Section~\ref{sec:linefull}.
The saved exact UNSAT results exclude every normalized pair outside the
two displayed planes over the unbounded real domain. Their weak-order
inequalities give precisely the parameter chambers above; the strict-order
version gives the strict lower bounds. Integer parameters are recovered
from $p=A_5-A_4,q=A_2$ in the first form and
$p=A_4-A_3,q=A_1+p$ in the second. The corresponding gap lists also prove
noncongruence. This is the reviewed Iw dependency, with its exact solver
trust retained, rather than a new solver run or a bounded-data inference.
\end{proof}

\begin{lemma}[BF: noncongruent real projection]\label{lem:bf}
Let $A,B$ be homometric six-element subsets of an abelian group $K$.
If some $\ell:K\to\R$ sends their six-atom lists to noncongruent
multisets, then, up to independent rigid motions and interchange,
$A,B$ have the group-valued Bloom formula
\[
 X=\{0,p,q-2p,2q-2p,2q,3q-p\},\qquad
 Y=\{0,p,q+2p,2q-p,2q+p,3q-p\},\quad p,q\in K.
\]
The conclusion depends on Iw and the finite integral certificate predicate
described in this proof.
\end{lemma}
\begin{proof}
Use Iw to normalize the two projected lists. The needed translations
subtract actual elements of $A$ or $B$, and reflections use inversion
and such translations, so they can be performed in $K$. Label equal
height atoms arbitrarily. The resulting projected coordinates have one
of~\eqref{eq:weighted-form-one}--\eqref{eq:weighted-form-two}. At this
stage their $p,q$ are real numbers; they are not assumed to lift to $K$.
The existence of group-valued parameters is what must be proved.

For each form compute the fifteen coefficient differences of its ordered
list. Both sides have the same fifteen distinct coefficient pairs. These
can be checked by subtracting entries of the displayed lists; as real
expressions the two common multisets are
\begin{align*}
 \mathcal D_1=\{&q-3p,2q-3p,q-2p,2q-2p,3q-2p,\\
 &q-p,2q-p,3q-p,q,2q,p,q+p,2q+p,2p,q+2p\},\\
 \mathcal D_2=\{&2q-3p,2q-2p,3q-2p,q-p,2q-p,3q-p,\\
 &q,2q,p,q+p,2q+p,2p-q,2p,q+2p,3p-q\}.
\end{align*}
Put $t=q/p$. Two unequal coefficient forms $up+vq$ and $u'p+v'q$
can agree only when $v\ne v'$ and
$t=(u'-u)/(v-v')$. Enumerate these rational slopes for the $\binom{15}{2}$
pairs and retain the appropriate chamber. This yields precisely $3,4,5$
for the first form and $3/2,5/3$ for the second. Away from these slopes
there is exactly one matching, the formal coefficient matching. All
forward real differences are positive there, so their signs are $+1$.

At an exceptional $t=a/b$, evaluate each coefficient pair at the integer
parameters $(p,q)=(b,a)$ and group edges into equal-height bins. In a bin
of size $k$, enumerate all $k!$ bijections. A zero bin also has $2^k$
independent sign choices; a positive bin has only the positive sign.
Thus the exact finite input is:
\begin{center}\small
\begin{tabular}{lrrrrr}
\toprule
Form & Generic & First slope & Second slope & Third slope & Total\\
\midrule
$1$ & $1$ & $128$ at $3$ & $16$ at $4$ & $4$ at $5$ & $149$\\
$2$ & $1$ & $128$ at $3/2$ & $4$ at $5/3$ & --- & $133$\\
\bottomrule
\end{tabular}
\end{center}
These $282$ signed matchings exhaust all possible lifts of every
noncongruent projected pair, because only the five explicitly computed
slopes permit additional edge coincidences. This is an argument over all
real parameters, not a parameter sample.

For clarity, the complete \emph{acceptance predicate} for this finite
step is as follows. Reconstruct every matrix from its signed labels by
\eqref{eq:formal-matching-row}. For each of the seven cases, compare the
saved labelled matchings with the full independent bin-permutation or
edge-by-edge enumeration, including signs on zero edges; require exact
set equality and absence of duplicates. Each matrix has an assigned
representative with both integer identities
\eqref{eq:two-way-lattice-cert}. Each representative has a certified
diagonalization~\eqref{eq:smith-contract}. Finally, in the resulting
exact group coordinates require one of:
\begin{enumerate}
 \item two distinct labels on the same endpoint have identical coordinates;
 \item explicit $p,q$, endpoint interchange if used, translations $t_X,t_Y$
 and a sign $\eta$ satisfy
 \[
 \operatorname{multiset}(E_i)=t_X+\operatorname{multiset}(X(p,q)),\quad
 \operatorname{multiset}(F_i)=t_Y+\eta\operatorname{multiset}(Y(p,q)),
 \]
 with the sides interchanged when specified. These equalities are checked
 by evaluating all twelve coefficient formulas; they are not inferred
 from a successful or failed parameter search.
\end{enumerate}
The saved package passes this predicate. Its $21$ distinct labelled
integral lattices consist of $8$ binary $\Z^2$ presentations, $3$ binary
$\Z$ presentations, $8$ binary $\Z\oplus C_2$ presentations, all with
explicit Bloom identities, and $2$ $\Z$ presentations with forced
collisions. Distinctness of the representative lattices can also be
certified by a row whose class is nonzero in the other's certified quotient;
it is ancillary to coverage.

An actual signed matching of the group pair supplies $G_M\to K$.
The collision outcomes contradict the six-element hypothesis. Every
remaining outcome has a Bloom identity already in $G_M$; applying the
homomorphism gives that identity in $K$, with distinctness inherited from
$A,B$. This proves BF, with precisely the Iw and finite-predicate
dependencies stated above.
\end{proof}

For example the apparently new torsion boundary configuration
\[
 \{0,g,3g,3g+h,7g,8g+h\}/
 \{0,2g+h,4g+h,7g,7g+h,8g+h\},\qquad 2h=0,
\]
is already Bloom: take $p=g+h,q=3g$ and reflect $Y,X$, respectively,
about $8g+h$. Torsion in a matching presentation alone does not prove
that its images are non-shadows.

\begin{corollary}[Intrinsic free-quotient dichotomy]\label{cor:free-dichotomy}
For a finitely generated universal presentation, either the endpoint
multisets in its full free quotient are congruent, or BF applies and the
pair is group-valued Bloom. In the congruent case a single labelled
alignment is valid for every real solution.
\end{corollary}
\begin{proof}
If no fixed vector congruence holds in the free quotient, the equalities for each choice of
sign, permutation and anchor can hold after a real functional only
on a proper subspace of the dual. There are finitely many proposals, so
Lemma~\ref{lem:finite-subspaces} supplies a functional outside all of
them. Its two projected multisets are noncongruent and BF applies.
In the other case use Lemma~\ref{lem:fixed-alignment}.
\end{proof}
% END WEIGHTED LINE INPUT AND BF

% BEGIN HR
\section{The high-free-rank cover and its direct \texorpdfstring{L3$^*$}{L3*} form}
\label{sec:formal-high-rank}
\begin{lemma}[HR]\label{lem:hr}
Let $M$ be an original signed matching presentation whose universal
endpoints have six distinct points and are $\TI$ distinct. If $d\ge3$,
then, conditional on Iw and the finite certificate predicate below,
\[
 G_M\cong C_2\oplus\Z^3.
\]
After relabellings and independent rigid motions its two configurations
are exactly
\begin{equation}\label{eq:hr-master}
 C=\{0,p,q,p-q+h\},\quad
 A=C\sqcup\{r,r+h\},\quad
 B=C\sqcup\{p-r,p-r+h\},\quad2h=0,
\end{equation}
where $p,q,r$ form a basis of the free quotient and $h$ generates $C_2$.
Every admissible specialization of this family is directly L3$^*$.
There are no $\TI$-distinct six-set universal presentations of rank
$d=4$ or $5$.
\end{lemma}
\begin{proof}[Reduction and exhaustive finite contract]
For each $v\in V_M$, Iw says that its real multisets are either congruent
or a labelled Bloom pair. There are finitely many label permutations,
orientations and choices of origin. Each fixed Bloom choice defines a
linear subspace of the anchored ten-coordinate space of dimension at
most two: all coordinates are fixed integer-linear combinations of
its two parameters, and anchoring removes the translations. Thus $V_M$
is covered by finitely many intersections with congruence subspaces and
Bloom subspaces. Since $\dim V_M\ge3$, none of the Bloom subspaces can
contain it. Lemma~\ref{lem:finite-subspaces} gives one fixed congruence
alignment. Apply Lemma~\ref{lem:fixed-alignment} and retain the original
integral presentation; only its real heights have $b_i=a_i$.

By~\eqref{eq:height-codimension}, the rational height span $S$ has
rank at most two. The primitive directions of all nonzero rows
\eqref{eq:height-rows}, with an overall sign removed, form the following
explicit $S_6$-invariant set $\mathcal D$:
\begin{align*}
 &(1,-1,0,0,0,0) &&\text{and its $15$ projective permutations},\\
 &(2,-1,-1,0,0,0) &&\text{and its $60$ projective permutations},\\
 &(1,1,-1,-1,0,0) &&\text{and its $45$ projective permutations}.
\end{align*}
To verify that these are all possibilities, subtract two signed
incidence roots. Equal supports give zero or a multiple of a root;
one shared vertex gives a root or the second shape; disjoint supports
give the third shape. The counts are respectively $\binom62$,
$6\binom52$, and $\binom64\binom42/2$. Hence $|\mathcal D|=120$.
Dividing by a coefficient gcd here does not divide any row of $M$.

Enumerate the zero span, every one-direction span, and every two-direction
span of $\mathcal D$, identifying rational spans by exact row reduction.
This enumeration is exhaustive because a span of rank at most two has
a basis chosen from its generating rows. Partition these spans by all
$720$ permutations of the six coordinates. A simultaneous permutation
of the aligned endpoints and reanchoring induces an invertible integral
change of the original generators, so one representative of each orbit
suffices. The saved enumeration has $4176$ spans in $25$ orbits.

For each representative $S$, supply an integer height vector
$z=(0,z_1,\ldots,z_5)$ with the certificate property
\begin{equation}\label{eq:generic-height-contract}
 v\cdot z=0\quad\Longleftrightarrow\quad v\in S,
 \qquad v\in\mathcal D.
\end{equation}
Such a vector exists: each $v\notin S$ cuts a proper hyperplane in the
anchored rational kernel of $S$; Lemma~\ref{lem:finite-subspaces} over
$\mathbb Q$ gives a point outside all of them, and clearing denominators
makes it integral. The stored vector must be checked against all $120$
directions. Its purpose is to record all forced signed-edge equalities,
not to approximate an arbitrary real solution.

Enumerate every permutation between the fifteen edge occurrences of
$z$ and another copy of $z$ that preserves absolute edge differences.
For a positive-length bin the sign is uniquely determined by the two
oriented differences. For a zero bin both signs are included independently.
If a bin has size $k$, it contributes $k!$, with an additional factor
$2^k$ only at zero. Reconstruct the full integral matrix for every
matching by~\eqref{eq:formal-matching-row}. Every original presentation
in this branch occurs in this collection after the justified changes of
generators: each of its matching rows restricts to a row in $S$, so
\eqref{eq:generic-height-contract} makes that signed equality true at
$z$. Conversely any equality at $z$ represents the same rational forced
height relation. Extra presentations with different integral rank may
be enumerated, which is harmless; none is omitted.

The acceptance predicate for the finite cover has four parts:
\begin{enumerate}
 \item Independently reconstruct $\mathcal D$ and all rank-at-most-two
 spans; check that the $25$ orbit lists are disjoint and their union is
 the full $4176$-span set. Check~\eqref{eq:generic-height-contract}.
 \item For every orbit representative compare the entire saved labelled
 matching set with an independent edge-by-edge enumeration. Require
 exact set equality, no repeated record, and the bin-factorial count.
 \item Each discarded record must carry either $cM=e_i-e_j$ or
 $cM=f_i-f_j$ for distinct labels, or a permutation $\pi$, sign $\eta$
 and anchor $k$ with six identities
 \[
 c_iM=e_i-\eta f_{\pi(i)}+\eta f_k.
 \]
 Check that $\pi$ is a bijection and all coefficients are integers.
 These identities prove collision or rigid equivalence in every image.
 \item Every remaining record has matrices satisfying
 \eqref{eq:smith-contract} with diagonal
 $(1,1,1,1,1,1,2,0,0,0)$. In the exact quotient coordinates, evaluate
 a supplied origin, rigid alignment and $p,q,r,h$ and require all twelve
 master coordinates~\eqref{eq:hr-master} to equal the universal endpoints.
 Require $h\ne0$, $2h=0$, and determinant $\pm1$ for the $3\times3$
 free-coordinate matrix of $p,q,r$.
\end{enumerate}
The last determinant and the order of $h$ prove that these parameters
generate the whole $C_2\oplus\Z^3$, rather than a proper sublattice:
their free components give a basis, and $h$ generates the remaining
torsion. Thus the record certifies the universal family itself.

The saved files pass this predicate, with $14475$ total matchings:
$13916$ collision records, $555$ universal rigid-equivalence records,
and four master-family records. The discard records contain $17246$
integer row identities. All four survivors have the same height orbit,
with span generated by
\[
 z_0-z_1-z_2+z_5=0,\qquad z_3-z_4=0;
\]
the generic stored vector is $(0,8,1,3,3,9)$ and that orbit has $128$
matchings. Independent replay checks every record; the numerical totals
alone are not the acceptance predicate. The preceding exhaustiveness
argument and these identities prove the structural part of HR. Since
the survivors have free rank three, none has rank four or five.
\end{proof}

\begin{proof}[Direct algebraic L3$^*$ identity]
In the group ring of any abelian target with $2h=0$, put
$\mathsf P=x^p$, $\mathsf Q=x^q$, $\mathsf R=x^r$, $\mathsf H=x^h$.
Then $\mathsf H^{-1}=\mathsf H$ and
\[
 C=1+\mathsf P+\mathsf Q+\mathsf P\mathsf Q^{-1}\mathsf H,
 \qquad D=\mathsf R(1+\mathsf H).
\]
Direct expansion gives
\begin{equation}\label{eq:hr-cosymmetric-defect}
 C-\mathsf PC^*=(1-\mathsf H)
                    (\mathsf Q-\mathsf P\mathsf Q^{-1}),
 \qquad (C-\mathsf PC^*)(1+\mathsf H)=0.
\end{equation}
Take the fixed block to be $C$, the moving dyad to be $D$, and
$s=p-2r$. Its cross distribution is $E=CD^*$, so
\begin{align*}
 x^{-s}E
 &=\mathsf P^{-1}\mathsf R^2 C\mathsf R^{-1}(1+\mathsf H)\\
 &=\mathsf R\mathsf P^{-1}C(1+\mathsf H)
  =\mathsf R C^*(1+\mathsf H)=E^*.
\end{align*}
This is exactly the L3$^*$ hypothesis. The shifted dyad is
\[
 x^sD=\mathsf P\mathsf R^{-1}(1+\mathsf H),
 \quad\text{with points }\{p-r,p-r+h\},
\]
so its endpoint is precisely $B$. Six-point distinctness supplies the
required disjointness. This proves homometry and the named direct move
without a global reflection or a choice of half of $s$.

For comparison the earlier L5 description is also valid: globally reflect
$B$ by $t\mapsto p-t$ to obtain $D+\mathsf PC^*$, with $D$ fixed and
$C$ moving. Equation~\eqref{eq:hr-cosymmetric-defect} gives
$DC^*=\mathsf P^{-1}DC$, the strict L5 first-cross-term condition.
The direct L3$^*$ description is the shorter form of the same family.
In a cyclic target six-point distinctness forces $h\ne0$, hence the
modulus is even and $h=n/2$. Its free parameters may otherwise specialize
arbitrarily, subject to endpoint distinctness and $\TI$ distinction.
\end{proof}

The BF and HR finite packages are respectively
\code{results/2026-09-30-six-bloom-cylinders.json} and
\code{results/2026-09-30-six-free-rank/}. The checker
\code{tests/test_six_cylinder_branches_review.py} independently rebuilds
the signed matching sets and verifies the integer identities, without
importing the builders or Smith/Hermite routines. The separate Iw sources,
formulas and solver evidence remain
\code{notes/2026-09-30-six-integer-theorem.md} and its logged adversarial
review. The present exposition does not claim new novelty or formal
proof-assistant verification.
% END HR
